\section{Methodology}

In this section, we present the overall architecture of our proposed Unified Dual-Mask Physical Model, termed GeometricDualMask, designed to address the persistent challenge of depth estimation on non-Lambertian surfaces in light field imaging. Different from previous works that heavily rely on the ideal Lambertian assumption, our framework explicitly decouples the physical properties of light field angular signals and adaptively routes distinct surface regions to specialized processing branches. As illustrated in Fig. 1, the architecture comprises a physics-driven decomposition module, a dual-branch feature extraction network, and a mask-guided fusion mechanism. This design effectively tackles the geometric violation caused by specular and translucent materials, enabling accurate and robust depth inference across complex real-world scenes.

The input to our model is a dense light field image consisting of $9 \times 9$ angular views, totaling 81 perspectives. Let $I(u,v)$ denote the local angular signal patch at spatial coordinates $(x,y)$, where $u$ and $v$ represent the angular coordinates ranging from 1 to 9. Before feeding into the network, the raw light field data undergoes a preprocessing pipeline. Specifically, we apply a tone-mapping enhancement  to improve the visibility of highlight and shadow regions, followed by spatial normalization to standardize the input intensity distribution. The preprocessed input tensor, with dimensions of $81 \times H \times W$, is then directed into the core physical decomposition module, where $H$ and $W$ denote the spatial height and width, respectively.

The overall data flow progresses through four primary modules to extract and integrate geometric features. First, the \textbf{Three-Layer Angular Signal Decomposition} module decouples the angular signal into physical components, formulated as:
\begin{equation}
\label{eq:eq7}
I(u,v) = I_{DC} + a \cdot u_c + b \cdot v_c + \varepsilon(u,v)
\end{equation}
Here, $I_{DC}$ represents the ambient illumination, $u_c$ and $v_c$ denote the centered angular coordinates, $a$ and $b$ are linear coefficients encoding disparity, and $\varepsilon(u,v)$ captures the material-dependent residual. Instead of employing frequency-domain analysis that fails under low-resolution angular sampling, we extract Geometric Pseudo-labels, such as the angular gradient, from the residual component to characterize material properties. Subsequently, the decomposed features are routed into two parallel branches. The \textbf{Lambertian EPI Branch} processes regions adhering to the Lambertian assumption by extracting epipolar plane image (EPI) slopes  and applying a continuous depth prior to regress the initial depth prediction $D_L$. Conversely, the \textbf{Non-Lambertian Geometric Branch} handles glossy and scattering surfaces. Recognizing that non-Lambertian reflections violate the linear epipolar constraint and form X-shaped patterns, this branch leverages bimodal distribution features and geometric convolutions to extract the depth prediction $D_{NL}$. Finally, the \textbf{Dual-Mask Generation & Fusion} module integrates these predictions. It generates pixel-wise masks, $M_L$ and $M_{NL}$, through a channel-balanced projection with a fusion dimension of 96. The final depth is computed via mask-weighted aggregation.

The output of the proposed architecture is a single-channel, high-precision depth map, denoted as $D_{final}$, with a target spatial resolution of $256 \times 256$. The final depth prediction is obtained by fusing the branch-specific outputs using the generated masks, expressed as:
\begin{equation}
\label{eq:eq8}
D_{final} = M_L \odot D_L + M_{NL} \odot D_{NL}
\end{equation}
where $\odot$ denotes the element-wise multiplication. To ensure the masks accurately separate Lambertian and non-Lambertian regions, we introduce a Mask BCE Loss during the training phase. This loss function supervises the mask generation using the geometric pseudo-labels derived from the angular gradient, formulated as $\mathcal{L}_{mask} = \text{BCE}(M, P_{angular})$, where $M$ is the predicted mask and $P_{angular}$ is the pseudo-label target. This explicit supervision forces the mask gap to reach the predefined threshold, guaranteeing that the dual-branch routing mechanism operates effectively and yields a globally consistent depth estimation.

\subsection{Three-Layer Angular Signal Decomposition (}

We design the Three-Layer Angular Signal Decomposition module to explicitly decouple the physical properties of light field angular signals, addressing the geometric violations caused by non-Lambertian surfaces. Previous works  primarily treat angular signals as holistic features or rely solely on epipolar plane image (EPI) (EPI) slopes, which often leads to ambiguous depth inference when specular or translucent materials disrupt the ideal Lambertian assumption. In contrast, our module decomposes the raw angular signal into distinct physical components. This physics-driven decoupling not only isolates the geometric depth cues from material-induced reflectance anomalies but also generates highly discriminative geometric pseudo-labels to supervise the subsequent dual-mask routing mechanism.

Let $I(u,v)$ denote the preprocessed local angular signal patch at a specific spatial coordinate, where $u$ and $v$ represent the angular coordinates ranging from 1 to 9. The input tensor for this module has a dimension of $81 \times H \times W$, representing the $9 \times 9$ angular views for each spatial pixel. To model the physical formation of the light field, we formulate the angular signal as a superposition of three distinct layers: the ambient illumination, the parallax-induced geometric shift, and the material-specific reflectance residue. Mathematically, this three-layer decomposition is expressed as:
\begin{equation}
\label{eq:eq9}
I(u,v) = I_{DC} + a \cdot u_c + b \cdot v_c + \varepsilon(u,v)
\end{equation}
where $I_{DC}$ represents the direct current component corresponding to the ambient illumination and base albedo. The variables $u_c = u - 5$ and $v_c = v - 5$ denote the centered angular coordinates, ensuring orthogonality in the subsequent regression. The coefficients $a$ and $b$ are the linear gradients along the horizontal and vertical angular dimensions, respectively, which encode the local parallax and thus the depth information. Finally, $\varepsilon(u,v)$ captures the residual component, representing high-frequency reflectance anomalies such as specular highlights and subsurface scattering that violate the linear parallax model.

To derive these components efficiently across the entire spatial domain, we vectorize the $9 \times 9$ angular patch into a column vector $\mathbf{i} \in \mathbb{R}^{81}$. The decomposition model is then reformulated in a matrix form:
\begin{equation}
\label{eq:eq10}
\mathbf{i} = \mathbf{X} \boldsymbol{\beta} + \boldsymbol{\varepsilon}
\end{equation}
where $\mathbf{X} = [\mathbf{1}, \mathbf{u}_c, \mathbf{v}_c] \in \mathbb{R}^{81 \times 3}$ is the design matrix constructed from the centered angular coordinates, $\boldsymbol{\beta} = [I_{DC}, a, b]^T \in \mathbb{R}^{3}$ is the parameter vector to be estimated, and $\boldsymbol{\varepsilon} \in \mathbb{R}^{81}$ is the residual vector. We employ ordinary least squares to estimate the parameter vector $\hat{\boldsymbol{\beta}}$, yielding:
\begin{equation}
\label{eq:eq11}
\hat{\boldsymbol{\beta}} = (\mathbf{X}^T \mathbf{X})^{-1} \mathbf{X}^T \mathbf{i}
\end{equation}
Given that the centered coordinates $\mathbf{u}_c$ and $\mathbf{v}_c$ are orthogonal and zero-mean, the matrix $\mathbf{X}^T \mathbf{X}$ is strictly diagonal. This property significantly reduces the computational complexity, allowing for rapid closed-form solutions during the forward pass. The estimated residual vector is subsequently computed as $\hat{\boldsymbol{\varepsilon}} = \mathbf{i} - \mathbf{X} \hat{\boldsymbol{\beta}}$.

Beyond extracting the physical components, this module derives high-discriminative geometric pseudo-labels to guide the dual-mask generation. We define the primary pseudo-label as the angular gradient, denoted as $\nabla_{ang}$, which quantifies the magnitude of the linear parallax shift. It is calculated as:
\begin{equation}
\label{eq:eq12}
\nabla_{ang} = \sqrt{a^2 + b^2}
\end{equation}
This angular gradient serves as a robust indicator of depth variations and geometric structures within the light field domain. Regions with high gradient magnitudes typically correspond to depth discontinuities or complex geometric boundaries, whereas flat regions exhibit near-zero gradients. By utilizing $\nabla_{ang}$ as a geometric pseudo-label, we provide explicit physical supervision for the mask generation module. This mechanism ensures that the network accurately distinguishes between Lambertian and non-Lambertian regions based on their underlying geometric consistency, thereby facilitating stable mask convergence and improving the overall depth estimation accuracy. With the methodological framework fully established, we now proceed to empirically validate its performance through extensive experiments.